\documentclass{article}
\usepackage{graphicx} % Required for inserting images
\usepackage{amsmath}
\usepackage{amsfonts}
\usepackage{amssymb}
\usepackage{amsthm}
\usepackage{tikz}
\usepackage[T1]{fontenc}
\usepackage[utf8]{inputenc}
\usepackage{mathtools}
\usepackage{tikz-cd}
\usepackage[only,llbracket,rrbracket]{stmaryrd}
\usepackage{manfnt}
%%%%%%%%%%%%%%%%%%%%
% CUSTOM DEFINITIONS
%%%%%%%%%%%%%%%%%%%%
\theoremstyle{definition}
\newtheorem{notation}{Notation}[section]

\theoremstyle{definition}
\newtheorem{definition}{Définition}[section]

\theoremstyle{remark}
\newtheorem{notabene}{NB}[section]

\newtheorem{theoreme}{Théorème}[section]
\newtheorem{lemme}[theoreme]{Lemme}
\newtheorem{corollaire}[theoreme]{Corollaire}
\newtheorem{proposition}[theoreme]{Proposition}
\newtheorem{exemple}[theoreme]{Exemple}
%%%%%%%%%%%%%%%%%%%%%%%%%%%%%%%%%%%%%%%%%%%%%%%%%%%%%%%%%%%%%%%%%%%
% IGNORER LES ERREURS DUES A [\Brackets] ELLES NE POSENT PAS 
% REELLEMENT PROBLEME
%%%%%%%%%%%%%%%%%%%%%%%%%%%%%%%%%%%%%%%%%%%%%%%%%%%%%%%%%%%%%%%%%%%
\DeclarePairedDelimiter\Brackets{\left\llbracket}{\right\rrbracket}

\title{Rapport TER}
\author{Lucas Rooy}
\date{Janvier 2026}

\begin{document}

\maketitle

\section{Introduction}

\section{Notations}
\begin{notation}[Composition]
    On utilise dans toute la suite l'ordre de composition catégorique. C'est-à-dire $f\circ g =x\mapsto g(f(x))$.
    De plus, pour simplifier les notations, on utilisera souvent la notation produit $f\circ g = f g$.
\end{notation}
\begin{notation}[Entiers]
On note dans toute la suite $\forall n\in \mathbb N, [n] = \{0\leq i < n, i\in\mathbb N\}$. De sorte que $\forall n \in \mathbb N, |[n]| = n$.
\end{notation}

\section{Présentation par générateurs et relations}
\subsection{D'un monoïde}

\begin{definition}[Monoïde engendré par générateurs et relations]
Soit $\Sigma$ un alphabet et $R \subseteq \Sigma^* \times \Sigma^*$ un ensemble de relations entre les mots de $\Sigma$.
On dit que $\Sigma^*/R$ est engendré pas les générateurs $\Sigma$ et les relations $R$.
\end{definition}

\begin{definition}[Monoïde présenté par générateurs et relations]
Soit $M$ un monoïde. Soit $\Sigma$ un alphabet, $R \subseteq \Sigma^* \times \Sigma^*$ et $f:\Sigma \to M$ une fonction \textit{d'interprétation}.
On étend $f$ au monoïde présenté (de façon unique) en $f^*:\Sigma^*/R \to M$.
On dit que $M$ est présenté par les générateurs $\Sigma$ et relations $R$ ssi. $f$ est bijective. 
\end{definition}

\subsection{D'une catégorie}
Dans le cas d'une catégorie, l'ensemble d'objets devient un graphe $G$.
\begin{definition}[Catégorie présentée par générateurs et relations]
cf. MacLane\cite{catswm}
\end{definition}

\section{La catégorie $\Delta$ et autres}
\subsection{Catégorie $\Delta$}
\subsubsection{Définition}

\begin{definition}[$\Delta$]
La catégorie $\Delta$ a pour objets $\{[n], n \in \mathbb N\}$, et pour hom-sets :

$$
\forall m, n \in \mathbb N, \text{Hom}_\Delta([m],[n]) = \{f, f\in \mathcal F([m], [n]) \text{ et }f\text{ est croissante}\}\cite{catswm}
$$
\end{definition}

\begin{notabene}
On peut représenter une flèche de $\text{Hom}_\Delta([m], [n])$ par un graphe. Exemple pour $f \in \text{Hom}_\Delta([3],[2])$ :

\begin{tikzpicture}
  \matrix (m) [matrix of nodes,
    row sep = 1em, column sep = 2em,
    nodes = {text depth = 1ex, text height = 2ex}
  ]
  {
    0 & 0 \\
    1 & 1 \\
    2 &   \\
  };
  \draw [-stealth]
    (m-1-1) edge (m-1-2)
    (m-2-1) edge (m-1-2)
    (m-3-1) edge (m-2-2)
  ;
\end{tikzpicture}
\end{notabene}

\subsubsection{Présentation}
\textit{Cette partie est largement recopiée du MacLane\cite{catswm}, car elle est utile pour la compréhension du rapport mais aussi car j'ai passé du temps dessus avant de passer à $\Theta$.}


On se donne un graphe G avec pour sommets $V = \{[n], n\in\mathbb N\}$.
Pour les arêtes, on introduit deux constructions :
\begin{itemize}
    \item Pour $n \in \mathbb N$ et $0 \leq i \leq n$, une fonction croissante injective $\delta_i^n:[n] \to[n+1]$ qui omet $i$ dans son image
    \item Pour $n \in \mathbb N$ et $0 \leq i \leq n$, une fonction croissante surjective $\sigma_i^n:[n+1]\to[n]$ qui envoie $i$ et $i+1$ sur la même image $i$.
\end{itemize}
Les arêtes sont alors ces $\delta_i^n$ et $\sigma_i^n$.

Dans la suite, les $n$ sont omis.

On soumet la catégorie libre engendrée par ce graphe $G^*$ aux relations :
\begin{itemize}
    \item $\delta_i\delta_j=\delta_{j+1}\delta_i$, $i \leq j$
    \item $\sigma_j\sigma_j=\sigma_i\sigma_{j+1}$, $i \leq j$
    \item $\sigma_j\delta_i=\delta_i\sigma_{j-1}$, $i < j$
    \item $\sigma_j\delta_i=1$, $i=j$, $i=j+1$
    \item $\sigma_j\delta_i=\delta_{i-1}\sigma_j$, $i>j+1$
\end{itemize}


\subsection{Catégorie $S$}
\subsubsection{Définition}
Il s'agit d'une restriction de $\Delta$ où les hom-sets sont : $\text{Hom}_S([n],[n])=S_n$ et $\forall m,n\in\mathbb N, m\neq n \implies \text{Hom}_\Delta([m], [n]) = \emptyset$. C'est-à-dire les bijections/permutations/symétries.
\subsubsection{Présentation}
On se donne un graphe G = (V,E) avec $V = \{[n], n\in\mathbb N\}$.
On introduit une construction pour les arêtes : $\forall n \in \mathbb N^*, \forall i \in [n-1], \tau_i^n:[n]\to[n]$ telle que :
\begin{equation}
   \tau_i^n(j) =
    \begin{cases}
      i+1 & j = i \\
      i & j = i +1 \\
      j & \text{sinon}
    \end{cases}       
\end{equation}
Autrement dit, $\tau_i$ est une permutation adjacente.
Ainsi, $E$ est donné par $\{\tau_i^n, n \in \mathbb N^*, i\in[n]\}$.

\begin{notabene}[Tri à bulles]
Par l'algorithme du tri à bulles\cite{knuthvol3}, toute permutation peut se mettre sous la forme d'un produit de permutations adjacentes, ce qui donne l'idée de cette présentation. 

Sans aller trop en détails, il s'agit d'un algorithme de tri déterministe par application de permutations adjacentes successives. Comme l'algorithme est déterministe, la suite de permutations adjacentes est canonique. On peut l'inverser pour obtenir le sens voulu.
\end{notabene}


Pour les règles :
\begin{itemize}
    \item $\tau_i\tau_i=id$
    \item $\tau_i\tau_j=\tau_j\tau_i$, $|i-j| > 1$
    \item $\tau_{i+1}\tau_i\tau_{i+1}=\tau_i\tau_{i+1}\tau_i$
\end{itemize}

\begin{figure}
    \centering
    \includegraphics[width=1\linewidth]{forme_normale_permutations.jpeg}
    \caption{Forme Canonique Permutations}
    \label{fig:enter-label}
\end{figure}

On a donc une forme canonique, elle est pratique, car la forme normale de l'identité est aussi l'identité.
Ainsi, la fonction d'interprétation étendue est bien bijective.

\subsection{Catégorie $F$}
\subsubsection{Définition}
Il s'agit d'une catégorie qui contient $\Delta$ où les hom-sets sont : $\text{Hom}_F([m],[n])=\mathcal F([m],[n])$.
\subsubsection{Présentation}
On crée encore $G=(V,E)$ avec $V =\{[n], n\in\mathbb N\}$.
Pour les arêtes, on prend l'union de celles utilisées pour $\Delta$ et pour $S$.
Les règles :
\begin{itemize}
    \item Les règles de la présentation de $\Delta$
    \item Les règles de la présentation de $S$
    \item ...
\end{itemize}
\section{La catégorie $\Theta$}
\subsection{Définition}
La catégorie $\Theta$ a été introduite dans un manuscrit de André Joyal\cite{thetajoyal}. Sa définition originelle est complexe et ne nous intéresse pas ici. Plutôt, on s'intéresse à la définition de Clemens Berger basée sur le \textit{wreath product} ou produit en couronne.

\begin{definition}[Wreath product ou \textit{Produit en couronne}]
cf. nLab.
\end{definition}

\begin{definition}[Catégories $\Theta_n$]
\[
  \forall n \in \mathbb N^*, \Theta_n =
    \begin{cases}
      \Delta & \text{si n = 1}\\
      \Delta\wr \Theta_{n-1} & \text{sinon}
    \end{cases}\cite{iteratedwp}
\]
\end{definition}
\begin{notabene}[Cas n=0]
En réalité, on peut commencer la définition à $n=0$ en utilisant la catégorie unitaire pour $\Theta_0$, ce qui aboutit à une seconde définition de $\Delta$.
\end{notabene}
\subsection{Présentation de la catégorie $\Theta$}
C'est une présentation \textbf{inductive}, basée sur la définition du produit en couronne.

Quelques notations :
\begin{itemize}
    \item $x$ désigne textuellement soit $\delta$ soit $\sigma$, pour garder le document concis,
    \item Pour tout $n\in\mathbb N^*$, $V_n$, les sommets du graphe de la présentation de $\Theta_n$,
    \item Pour $n\in\mathbb N$, $E_n$, les arêtes du graphe de la présentation de $\Theta_n$,
    \item Pour $n\in\mathbb N$, $\mathcal R_n \subseteq (E_n^*)²$ les relations de la présentation de $\Theta_n$,
    \item Pour une arête $y$ de $G_n$, son interprétation sera notée $\Brackets{y}_n$. C'est une fonction $\Brackets{.}_n : E_n\to \Theta_n$ compatible avec les relations $\mathcal R_n$. Elle possède une extension naturelle à la catégorie libre sur le graphe notée $\Brackets{.}_n^*$ (il s'agit alors d'un foncteur).
\end{itemize}

La présentation de $\Theta_1$ est déjà connue

Soit $n\in\mathbb N^*$.
On connaît $((V_n, E_n), \mathcal R_n)$ la présentation de $\Theta_n$. On va construire la présentation de $\Theta_{n+1}$.

\subsubsection{Objets}
Les objets de $\Theta_{n+1}$ sont donnés directement par la définition du produit en couronne, on pose :
\[
    V_{n+1} = \{(n_1,\dots,n_m)\in V_n^m,m\in\mathbb N,\}
\]

Leur interprétation est l'identité.

\subsubsection{Morphismes}
On introduit la notion de flèches «internes» et «externes».
\begin{itemize}
    \item Les flèches internes sont héritées de $\Theta_n$, elles agissent sur les objets présents dans le tuple objet.
    \item Les flèches externes travaillent sur la liste des objets directement, de la même façon que les flèches de $\Delta$.
\end{itemize}
Cette intuition sera confirmée par l'expression formelle de la fonction d'interprétation plus bas.

\textbf{Morphismes internes}
On les noter $\mathcal I_{n+1}$.
\[
\mathcal I_{n+1}=\{x^{(n_1,\dots,n_m)}_{i,j},m\in\mathbb N,(n_1,\dots,n_m)\in V_{n+1},i\in[m],j\in n_i\}
\]
Soit $x^{(n_1,\dots,n_m)}_{i,j}\in \mathcal I_{n+1}$. Alors $x^{n_i}_j \in \mathcal I_n$ existe nécessairement.
\[\
\Brackets{x^{(n_1,\dots,n_m)}_{i,j}}_{n+1}^* : (k_1,\dots,k_m) \mapsto (k_1,\dots,\Brackets{x^{n_i}_j}_n^*(k_i),\dots,k_m)
\]

\textbf{Morphismes externes}
On peut les noter $\mathcal E_{n+1}$
\[
\mathcal E_{n+1}=\{x^{(n_1,\dots,n_m)}_i,m\in\mathbb N,(n_1,\dots,n_m)\in V_{n+1},i\in [m]\}
\]

Soit $x^{(n_1,\dots,n_m)}_i \in \mathcal E_{n+1}$.
\begin{itemize}
    \item Si $x=\delta$.\[
\Brackets{x^{(n_1,\dots,n_m)}_i}_{n+1}^* : (n_1,\dots,n_m)\mapsto (n_1,\dots,\hat{n_i},\dots,n_m)
\]
\item Si $x=\sigma$. \[
\Brackets{x^{(n_1,\dots,n_m)}_i}_{n+1}^* : (n_1,\dots,n_m)\mapsto (n_1,\dots,n_i,n_i,n_{i+2},\dots,n_m)
\]
\end{itemize}

\subsubsection{Relations}
\textbf{Les relations héritées}
Soient
\[
\begin{split}
u & = x^{(n_{1,1},\dots,n_{1,m_1})}_{i_1,j_1}\dots x^{(n_{k,1},\dots,n_{k,m_k})}_{i_k,j_k}\\
v & = x^{(n'_{1,1},\dots,n'_{1,m'_1})}_{i'_1,j'_1}\dots x^{(n'_{k',1},\dots,n'_{k',{m'}_{k'}})}_{{i'}_{k'},{j'}_{k'}}
\end{split}
\]

On se donne aussi,

\[\begin{split}
u' & = x^{n_{1,i_1}}_{j_1}\dots x^{n_{k,i_k}}_{j_k}\\
v' & = x^{n'_{1,i'_1}}_{j'_1}\dots x^{n'_{k',{i'}_{k'}}}_{{j'}_{k'}}
\end{split}
\]

Alors, $(u,v)\in \mathcal R_{n+1}$ ssi. $(u',v')\in\mathcal R_n$.
Cette classe de relations permet de propager les relations de $\Theta_n$.

\textbf{Les relations simpliciales sur les flèches externes}
Soient
\[
\begin{split}
    u & = x^{n_{1,1},\dots,n_{1,m_1}}_{i_1}\dots,x^{(n_{k,1},\dots,n_{k,m_k})}_{i_k}\\
    v & = x^{n'_{1,1},\dots,n'_{1,m'_1}}_{i'_1}\dots,x^{(n'_{k',1},\dots,n'_{k',{m'}_{k'}})}_{{i'}_{k'}}
\end{split}
\]

On se donne aussi,

\[
\begin{split}
    u' & = x^{m_1}_{i_1}\dots,x^{m_k}_{i_k}\\
    v' & = x^{m'_1}_{i'_1}\dots,x^{m'_{k'}}_{{i'}_{k'}}
\end{split}
\]

Alors, $(u,v)\in\mathcal R_{n+1}$ ssi. $(u',v')\in\mathcal R_\Delta$.

\textbf{Echange loin des internes}
Lorsque des flèches internes ne concernent pas les mêmes cases, on peut les échanger.
Soient $m\in\mathbb N$, $(n_1,\dots,n_m) \in V_{n+1}$, $i,i'\in[m]$ tq $i\neq i'$, $(j,j')\in(n_i,n_{i'})$.
Alors, $(x^{(n_1,\dots,n_m)}_{i,j}x^{(n_1,\dots,n_m)}_{i',j'},x^{(n_1,\dots,n_m)}_{i',j'}x^{(n_1,\dots,n_m)}_{i,j}) \in \mathcal R_{n+1}$

\textbf{Relations croisées}
Lorsqu'une flèche interne suit une flèche externe, on peut la faire passer «avant» la flèche externe, sous condition de modifier les indices proprement.

Pour mieux comprendre les cas à considérer, deux dessins :

\begin{tikzcd}
0 \arrow[dd] & 1 \arrow[dd] & 2 \arrow[rdd] & 3 \arrow[rdd] & 4 \arrow[rdd] &   \\
              &               &                &                &                &   \\
0             & 1             & 2              & 3              & 4              & 5
\end{tikzcd}

Les deux premiers points correspondent à un intervalle de \textbf{commutation}.

Le suivant correspond à un intervalle de \textbf{disparation}.

Les trois derniers correspondent à un intervalle de \textbf{décalage}.

\begin{tikzcd}
0 \arrow[dd] & 1 \arrow[dd] & 2 \arrow[dd] & 3 \arrow[ldd] & 4 \arrow[ldd] & 5 \arrow[ldd] \\
              &               &               &                &                &                \\
0             & 1             & 2             & 3              & 4              &               
\end{tikzcd}

Les deux premiers points correspondent à un intervalle de \textbf{commutation}.

Le suivant correspond à un intervalle de \textbf{dédoublement}.

Les deux derniers correspondent à un intervalle de \textbf{décalage}.

\[
  \delta^{(n_1,\dots,n_m)}_i x^{(n_1,\dots,\hat{n_i},\dots,n_m)}_{k,j} =
    \begin{cases}
      x^{(n_1,\dots,n_m)}_{k,j} \delta^{(n_1,\dots,n_m)}_i& \text{si $k<i$}\\
      \delta^{(n_1,\dots,n_m)}_i & \text{si $i=k$}\\
      x^{(n_1,\dots,n_m)}_{k-1,j} \delta^{(n_1,\dots,n_m)}_{i} & \text{si $k>i$}
    \end{cases}
\]
\[
  \sigma^{(n_1,\dots,n_m)}_i x^{(n_1,\dots,\hat{n_i},\dots,n_m)}_{k,j} =
    \begin{cases}
      x^{(n_1,\dots,n_m)}_{k,j} \sigma^{(n_1,\dots,n_m)}_i& \text{si $k<i$}\\
      x^{(n_1,\dots,n_m)}_{k,j}x^{(n_1,\dots,n_m)}_{k+1,j}\delta^{(n_1,\dots,n_m)}_i & \text{si $i=k$}\\
      x^{(n_1,\dots,n_m)}_{k+1,j} \sigma^{(n_1,\dots,n_m)}_{i} & \text{si $k>i$}
    \end{cases}
\]

\subsubsection{Forme canonique unique}
On démontre par récurrence à droite que tout mot de composition de flèches de $G_{n+1}$ possède une forme canonique unique de la forme $\text{internes}\circ\text{externes}$.

La démonstration n'est pas très intéressante, car les cas difficiles ont déjà été faits dans la présentation de $\Delta$. En effet, les relations donnent directement le passage des internes le plus à gauche possible.

\textdbend Pour assurer l'unicité, il faut néanmoins faire attention et imposer un ordre sur la commutation des internes loins (par exemple, l'ordre croissant sur le premier indice (sélection de la case)), ainsi que sur la commutation des externes (comme fait pour présenter $\Delta$). 

\textbf{A RETRAVAILLER, MAIS IDEE PREUVE :}

Mots $\text{internes} \circ \text{externes}$, ajout d'un morphisme à droite.
\begin{itemize}
\item Si il est externe : les relations le mettent dans le bon ordre. Avec la relation d'ordonnancement: on a un ordre. OK.
\item Si il est interne : les relations la font percoler jusqu'aux autres internes. Avec les relations de commutation : premier ordre. Parmi celles qui ne commutent pas, on applique bien les relations héritées pour obtenir une forme normale
\end{itemize}

La mise sous forme canonique des deux «côtés» est ensuite directe par la présentation de $\Delta$.

Pour un morphisme de $G_{n+1}^*$ $f$, on pourra noter $\overline{f}$ sa classe d'équivalence pour les relations $\mathcal R_{n+1}$, et $\tilde{\overline{f}}$ la forme canonique de sa classe d'équivalence.

\subsubsection{Présentation}

On a donc :

\begin{itemize}
    \item Une bijection entre les morphismes de la catégorie libre $G_{n+1}^*$ et un ensemble de formes canoniques.
    \item Une fonction d'interprétation $\Brackets{.}_{n+1}^*:G_{n+1}^*\to \Theta_{n+1}$. On cherche à montrer sa bijection modulo les classes d'équivalences induites par $\mathcal R_{n+1}$, donc, on peut ne s'intéresser qu'à la bijection $\Brackets{.}_{n+1} : G_{n+1}^*/\mathcal R_{n+1} \to \Theta_{n+1}$. En l'occurrence, on utilisera comme représentants de chaque classe d'équivalence, la forme canonique associée.
\end{itemize}

Je fais la preuve de bijection par injection + surjection du foncteur.

\begin{itemize}
    \item \textbf{Injectivité}. On a forcément $\tilde{\overline{\text{Id}}}=\text{Id}$. Donc, $\Brackets{\tilde{\overline{\text{Id}}}}_{n+1}^*=\Brackets{\text{Id}}_{n+1}^*=\text{Id}$. C'est donc bien injectif.
    \item \textbf{Surjectivité}. C'est le cas difficile. Il faut donner la forme canonique de n'importe quel morphisme de $\Theta_{n+1}$. D'après Berger\cite{iteratedwp} (Remark 3.15, preprint), un morphisme de $\Theta_{n+1}$ est vu comme le produit d'un morphisme interne et d'un morphisme externe. C'est donc bien la construction voulue et l'interprétation étoilée est bien surjective. \textbf{MODIFIER CE PARAGRAPHE POUR REFAIRE LA PREUVE. JE NE SUIS PAS SÛR DE CELLE DE BERGER (SUPPR. DANS L'ARTICLE SOUMIS)}
\end{itemize}

Ainsi, on obtient bien une présentation par générateurs et relations de toutes les $\Theta_n$ catégories.

\begin{thebibliography}{9}
\bibitem{iteratedwp}
Clemens Berger (2005) \emph{Iterated wreath product of the simplex category and
iterated loop spaces}, Advances in Mathematics.

\bibitem{thetajoyal}
Andre Joyal (1997) \emph{Disks, duality and $\Theta$-categories}, draft.

\bibitem{knuthvol3}
Donald E. Knuth (1998) \emph{The Art of Computer Programming Vol. 3}, Stanford Computer Science

\bibitem{catswm}
Saunders Mac Lane (1978) \emph{Categories for the Working Mathematician}, Springer Nature.
\end{thebibliography}
\end{document}

